We study what is shared, and what is optimizer-dependent, in teacher-forced language-model loss.  The recent Relative Generalization Invariance (RGI) result reports that token-wise loss differences can remain highly similar across training choices, but it does not identify the finite-update terms responsible for a Transformer trajectory.  We decompose an arriving-batch loss into a frozen-data difficulty and a local response.  For an update from \(\theta_0\) to \(\theta_t\), the response has an exact linear-work term \(A=\nabla_\theta\ell(\theta_0)^\top(\theta_t-\theta_0)\), a nonnegative logit-space curvature term \(Q=\operatorname{KL}(p_0\Vert p_t)\), and a representation residual.  Its second-order logit approximation is \(A+Q_2\), with \(Q_2=\tfrac12\operatorname{Var}_{p_0}(\Delta z)\).

Across three permutations of a fixed FineWeb pool and SGD or Adam continuation of Pythia-70M, a common step-32,000 reference explains 96.657--97.582\% of arriving-loss variance; each algorithm's frozen start explains 99.940--99.997\% in the registered near and far windows.  On the residual slow signal \(L-D\), observed-displacement \(A+Q_2\) has 4.97--9.73\% relative RMS error for SGD and 1.19--2.82\% for Adam, with at least 96.87\% variance explained.  Signed covariance shows negative linear work and positive curvature cancel, and both terms change with the optimizer.  A single scalar probe does not close the batch-specific response, and the optimizer difference fails the 10\% response gate.  A fixed-teacher continuation closes the same local decomposition, but its cross-prefix optimizer gap still fails a strict constant-offset test.  We therefore obtain a quantitative local finite-displacement mechanism and a clear boundary: the evidence does not establish a universal strong-RGI mechanism for natural language models.
